% =============================================================================
% TABLE 1 -- Before/after comparison  (Reviewer comment #4)
% =============================================================================
\begin{table*}[!htbp]
  \centering
  \caption{Summary of operational and sustainability outcomes before and
           after the Kirkpatrick-based intervention.}
  \label{tab:before-after}
  \small
  \begin{tabular}{@{}p{6.5cm}rrr@{}}
    \toprule
    \textbf{Metric} & \textbf{Before} & \textbf{After} & \textbf{Variation} \\
    \midrule
    Training failures (rate, \%)               & 100        & 72         & $-28\%$ \\
    Regulatory non-conformances (n / audit)    & 100        & 65         & $-35\%$ \\
    First-attempt pass rate (Level~2, \%)      & 78         & 86         & $+8$ pp \\
    Work-card consultation (Level~3, \%)       & 71         & 94         & $+23$ pp \\
    Quality costs (USD M / year)               & $\approx 4.2$ & $\approx 3.0$ & $\approx -1.2$ \\
    Turnaround time (index, baseline $=100$)   & 100        & 94         & $-6\%$ \\
    Customer satisfaction (0--5 scale, mean) & n.a.\textsuperscript{a} & 4.7 & --- \\
    \bottomrule
  \end{tabular}\\[2pt]
  \footnotesize
  \emph{Notes.} Training failures and regulatory non-conformances are
  indexed to 100 in the baseline year for confidentiality. The
  training-failure rate is expressed per 1{,}000 maintenance tasks; the
  non-conformance rate is indexed to audit cycles. Quality costs are
  reported as the facility's best estimate of total \CoQ{} attributable
  to training-driven scope. Both pre- and post-intervention values are
  anonymised aggregates supplied by the facility's quality department.
  \textsuperscript{a}~The baseline year did not employ the structured
  customer-survey instrument used in the intervention year; a directly
  comparable pre-intervention mean is therefore not available.
\end{table*}
